\section{Introducción}
    Mejorar la estética de una imagen no es la única forma para facilitar la extracción 
    de patrones, estructuras u otros elementos de la misma, también existen las 
    transformaciones afines que van a permitir modificar la perspectiva  
    de una imagen con el propósito de alinear los elementos de interés con los 
    ejes o ampliarlos para mejorar su detección. Y como se ha mencionado en 
    anteriores proyectos, estas transformaciones pueden ser parte del 
    preprocesamiento durante el entrenamiento de una CNN para expandir y generalizar 
    el aprendizaje del modelo hacia escenarios donde las imágenes no son exactamente 
    perfectas o ideales.\\

    Para la realización de este proyecto, se aplican transformaciones afines tanto unitarias 
    como compuestas para comparar los efectos que estas tienen sobre diferentes composiciones 
    de elementos. En general, cada transformación afín trae consigo un potencial beneficio 
    para muchas imágenes que no se encuentran en sus escenarios ideales, por ello, comprender 
    cómo modifican una imagen facilitará su uso en posteriores proyectos.